% Includable lemma and proof. Independently checked; no expansion of the
% fixed exclusion product is claimed.
\section{Coefficient bounds for a fixed elliptic exclusion product}
\label{sec:fixed-anchor-height}

The finite size condition in the shared-point construction uses fixed
multiples of a fixed elliptic point.  The number of factors is extremely
large, but their coefficients admit elementary uniform bounds.  This
section gives an exact algorithm defining those coefficients and proves
bounds without computing the full product.

Put $k=\mathbb Q(\sqrt2)$.  Every $x\in k$ has a unique expression
\[
 x=\frac{a+b\sqrt2}{c},\qquad
 a,b,c\in\mathbb Z,\quad c>0,\quad \gcd(a,b,c)=1.
\]
We call $(a,b,c)$ its primitive common-denominator triple and set
\[
 h_0(x)=\left\lceil\log_2\max(1,|a|,|b|,c)\right\rceil.
\]
Thus the triple of zero is $(0,0,1)$ and $h_0(0)=0$.  This elementary
coefficient height is used only in this section; it is not the canonical
elliptic height or the absolute logarithmic Weil height.

\begin{lemma}\label{lem:fixed-triple-height}
For $x,y\in k$ one has
\[
 \begin{aligned}
 h_0(x+y)&\le h_0(x)+h_0(y)+1,\\
 h_0(xy)&\le h_0(x)+h_0(y)+2,
 \end{aligned}
\]
and, for $x\ne0$,
\[
                         h_0(x^{-1})\le2h_0(x)+2.
\]
\end{lemma}
\begin{proof}
Write the primitive triples of $x,y$ as $(a,b,c)$ and
$(a',b',c')$.  Before cancellation their sum has triple
\[
 (ac'+a'c,\,bc'+b'c,\,cc')
\]
and their product has triple
\[
 (aa'+2bb',\,ab'+a'b,\,cc').
\]
The coordinates are bounded respectively by
$2^{h_0(x)+h_0(y)+1}$ and $3\cdot2^{h_0(x)+h_0(y)}$.
Passing to a primitive triple can only decrease their maximum.  For
the inverse use
\[
                  x^{-1}=\frac{ca-cb\sqrt2}{a^2-2b^2}.
\]
The denominator is nonzero because $x\ne0$.  Its absolute value is
at most $3\cdot2^{2h_0(x)}$, and the numerator coefficients are at
most $2^{2h_0(x)}$.  Changing all signs if necessary and cancelling
the common divisor proves the last bound.
\end{proof}

Consider the fixed elliptic curve and point
\[
 \begin{gathered}
 E:\quad y^2=x^3-D_E x,\qquad D_E=11-6\sqrt2,\\
 P=\left(\frac{57-28\sqrt2}{8},
          \frac{-660+417\sqrt2}{32}\right).
 \end{gathered}
\]
Direct substitution verifies $P\in E(k)$.  For a finite point
$R=(x,y)$ put $H_0(R)=\max(h_0(x),h_0(y))$, and set
$H_0(\mathcal O)=0$ for the identity.  Here $h_0(D_E)=4$,
$h_0(x(P))=6$, $h_0(y(P))=10$, and $H_0(P)=10$.

\begin{lemma}\label{lem:fixed-point-binary-height}
For every integer $n\ge1$, the point $nP$ is defined by exact
arithmetic in $k$ and
\[
                         H_0(nP)\le17n^8.
\]
Consequently, if $nP\ne\mathcal O$ and
$x(nP)=(a_n+b_n\sqrt2)/c_n$ is written in its primitive
common-denominator triple, then
\[
               \max(|a_n|,|b_n|,c_n)\le2^{17n^8}.
\]
No nontorsion assumption is needed for these inequalities.
\end{lemma}
\begin{proof}
We first establish bounds for doubling and addition of the fixed point.
Let $R=(x,y)$ be finite and $H=H_0(R)$.  If $y\ne0$, the usual
doubling formulas are
\[
 m=\frac{3x^2-D_E}{2y},\qquad
 x_2=m^2-2x,\qquad y_2=m(x-x_2)-y.
\]
Lemma~\ref{lem:fixed-triple-height}, together with
$h_0(2)=1$ and $h_0(3)=2$, gives successively
\[
 \begin{array}{c|c}
 \text{quantity}&\text{bound for its coefficient height}\\ \hline
 3x^2-D_E&2H+11\\
 (2y)^{-1}&2H+8\\
 m&4H+21\\
 x_2&9H+48\\
 x-x_2&10H+49\\
 y_2&15H+73.
 \end{array}
\]
Thus $H_0(2R)\le15H+80$.  When $R=\mathcal O$ or $y=0$,
the result is the identity, so the same bound holds using its
specified height zero.

For $R\ne\mathcal O,P,-P$, write $P=(x_P,y_P)$ and use
\[
 m=\frac{y-y_P}{x-x_P},\qquad
 x_+=m^2-x-x_P,\qquad y_+=m(x-x_+)-y.
\]
The denominator is nonzero: points of a short Weierstrass curve with
the same abscissa are opposite.  Using the uniform bounds
$h_0(x_P),h_0(y_P)\le10$ gives
\[
 \begin{array}{c|c}
 \text{quantity}&\text{bound for its coefficient height}\\ \hline
 y-y_P,\ x-x_P&H+11\\
 (x-x_P)^{-1}&2H+24\\
 m&3H+37\\
 x_+&7H+88\\
 x-x_+&8H+89\\
 y_+&12H+129.
 \end{array}
\]
It follows that $H_0(R+P)\le12H+129$.  The exceptional cases
obey this bound as well.  If $R=\mathcal O$ the output has height
10, and if $R=-P$ it has height zero.  If $R=P$, the doubling
estimate above with $H_0(P)=10$ bounds $H_0(2P)$ by 223,
which is smaller than $12\cdot10+129$.

To compute $nP$, begin with the leading binary digit of $n$, so the
initial point is $P$.  Read each subsequent digit from left to right,
doubling and then, for a digit one, adding $P$.  Every step is a
well-defined group operation; use the identity and the exceptional
cases just specified instead of a formula with zero denominator.
One binary step replaces a height bound $H$ by at most
\[
                     12(15H+80)+129=180H+1089.
\]
The same upper bound covers a digit zero.  Therefore the next height
$H'$ satisfies $H'+7\le180(H+7)$.  There are
$r=\lfloor\log_2 n\rfloor$ digits after the leading digit, and hence
\[
 H_0(nP)+7\le17\cdot180^r
       \le17\cdot(2^r)^8\le17n^8,
\]
since $180<2^8$.  This proves the asserted weaker bound
$H_0(nP)\le17n^8$ and the bound for the primitive triple.
\end{proof}

\subsection{A bound in a finite algebra}

Let $R=\mathbb Z[X_1,\ldots,X_d]$ and let $A$ be a commutative
$R$-algebra free of positive finite rank, with a specified basis.  For
$a\in A$ let $M_a$ be its multiplication matrix.  If
$\|f\|_1$ denotes the sum of the absolute values of the coefficients
of a polynomial $f\in R$, define
\[
              \|M\|=\max_j\sum_i\|M_{ij}\|_1.
\]
This norm satisfies the triangle inequality, is submultiplicative,
and gives $\|M_c\|=|c|$ for an integer scalar $c$.
The same arguments apply over rational coefficients with the
corresponding coefficient norm.

Suppose $N_0,N_1,D_0\in A$ and, for a real number $b\ge0$,
\[
        \|M_{N_0}\|,\ \|M_{N_1}\|,\ \|M_{D_0}\|\le2^b.
\]
For each $n\ge1$ with $nP\ne\mathcal O$, define
\[
 \mathcal E_n=(c_nN_0-a_nD_0)^2
                       -2(c_nN_1-b_nD_0)^2\in A;
\]
when $nP=\mathcal O$, define $\mathcal E_n=1$.  For an integer
$M\ge1$ put $\mathcal E_{\le M}=\prod_{n=1}^M\mathcal E_n$.

\begin{proposition}\label{prop:fixed-exclusion-norm}
The multiplication matrices of these factors satisfy
\[
 \begin{aligned}
  \|M_{\mathcal E_n}\|&\le2^{34n^8+2b+4},\\
  \|M_{\mathcal E_{\le M}}\|&\le2^{(2b+38)M^9}.
 \end{aligned}
\]
The fixed coefficients in the product are specified by the binary
algorithm of Lemma~\ref{lem:fixed-point-binary-height}; they are not
additional quantified variables.
\end{proposition}
\begin{proof}
For a finite $nP$, let $h_n=h_0(x(nP))$.  The triangle inequality
gives
\[
 \|M_{c_nN_0-a_nD_0}\|,\ \|M_{c_nN_1-b_nD_0}\|
       \le2^{h_n+b+1}.
\]
Submultiplicativity and the factor's two terms now give
\[
 \|M_{\mathcal E_n}\|
     \le3\cdot2^{2h_n+2b+2}
     \le2^{34n^8+2b+4}.
\]
The identity factor when $nP=\mathcal O$ obeys the same bound.
Multiplying and using $\sum_{n=1}^M n^8\le M^9$ and $M\le M^9$
gives
\[
 \|M_{\mathcal E_{\le M}}\|
 \le2^{34\sum_{n=1}^M n^8+(2b+4)M}
 \le2^{(2b+38)M^9},
\]
as claimed.  This reasoning does not require expansion of the product.
\end{proof}

At a rational specialization for which $D_0\ne0$, the value
$\mathcal E_n$ for a finite $nP$ is zero precisely when
\[
               \frac{N_0+N_1\sqrt2}{D_0}=x(nP).
\]
Indeed $u^2-2v^2=0$ with $u,v\in\mathbb Q$ forces $u=v=0$.
If the left side is the abscissa of a finite point $Q\in E(k)$,
the equality is equivalent to $Q=\pm nP$.  Thus the product,
together with the existing denominator guard, excludes exactly these
finitely many abscissas.  The convention for $nP=\mathcal O$ makes
the indexed algorithm total and introduces no unwanted exclusion.
